\documentclass[11pt]{article}
\usepackage[a4paper,margin=25mm]{geometry}
\usepackage{amsmath,amssymb,amsthm,booktabs}
\newtheorem{theorem}{Theorem}
\newtheorem{lemma}[theorem]{Lemma}
\newcommand{\Aut}{\operatorname{Aut}}
\title{Primitive residues in rank two}
\author{Human Compression project}
\date{Working proof, 14 October 2026}
\begin{document}
\maketitle

Let $G$ be finite. We use right cosets, with incidence given by nonempty
intersection. Firmness means that every panel lies in at least two chambers.
For a rank-two geometry, residual weak primitivity requires a primitive
$G$-action on at least one type and primitive actions of both vertex
stabilizers on their respective neighbours. Thus it includes conditions
on the rank-one residues. No double transitivity is assumed.

\begin{theorem}\label{t:complete}
All connected, firm, flag-transitive, residually weakly primitive
rank-two geometries for $G$ satisfying the intersection property are
obtained by the following construction.
\begin{enumerate}
\item Choose one representative $A$ of each conjugacy class of maximal
proper subgroups of $G$.
\item Choose one representative $H$ of each $A$-conjugacy class of maximal
proper subgroups of $A$.
\item Take every actual subgroup $B$ with $H<B<G$, retaining it precisely
when $A\cap B=H$ and $H$ is maximal in $B$.
\item Retain the pairs passing Lemma~\ref{l:allpairs}. Identify pairs under
automorphisms of $G$ and interchange of their entries.
\end{enumerate}
The final identification classifies the specified $G$-actions. Distinct
resulting classes may require a further comparison if one wishes to forget
the specified action and classify abstract incidence geometries.
\end{theorem}
\begin{proof}
Choose an incident pair in such a geometry. Its stabilizers $A,B$ identify
the geometry with the incidence structure of their cosets, by flag
transitivity. Interchange the types so that $A$ is maximal in $G$, and
conjugate to the chosen representative. The common stabilizer $H=A\cap B$
is maximal in both $A$ and $B$: in a transitive permutation action,
primitivity is equivalent to maximality of a point stabilizer. Conjugating
within $A$ now puts $H$ in the second list. The resulting $B$ is one of
the actual overgroups required in the third step. This proves completeness
before the incidence test and equivalence comparison.

Conversely, the two proper containments $H<A,B$ give firmness. Since
$H<B$ and $A\cap B=H$, the subgroup $B$ is not contained in $A$.
Maximality of $A$ therefore gives $\langle A,B\rangle=G$, which proves
connectivity. Every incident pair of cosets contains a common element,
so it is a translate of $(A,B)$; hence the action is flag-transitive.
The three maximality conditions give precisely the primitive actions
required above. These are all the nonempty residues that need checking
in rank two. Lemma~\ref{l:allpairs} supplies the intersection property.
An automorphism carrying one pair to another, possibly interchanging
its entries, induces the corresponding incidence isomorphism.
\end{proof}

\begin{lemma}\label{l:allpairs}
For a pair obtained in the third step, let $T_A,T_B$ be right transversals
of $H$ in $A,B$. For each $a\in T_A\setminus H$, put
\[
 c(a)=\bigl|\{b\in T_B:\text{ some }d\in T_B
                              \text{ has }b(da)^{-1}\in A\}\bigr|.
\]
The intersection property holds if and only if either $c(a)=1$ for every
such $a$, or \mbox{$[B:H]=[G:A]$}.
\end{lemma}
\begin{proof}
The neighbours of $B$ are $Ab$ for $b\in T_B$, while those of $Ba$ are
$Ada$ for $d\in T_B$. Both lists have distinct entries because
$A\cap B=H$. Equality of these two right cosets is exactly the displayed
membership test, so $c(a)$ counts their common neighbours. The cosets
$Ba$ with $a\in T_A$ exhaust the neighbours of $A$. The group $A$ is
transitive on these neighbours; consequently every pair of distinct
neighbours of $A$ can be moved to $(B,Ba)$ for one of the tested cosets.
Transitivity on vertices of the first type now shows that the condition
$c(a)=1$ for all $a$ is equivalent to having no quadrangle anywhere in
the incidence graph. This is the partial linear space condition.

The other allowed possibility is a generalized digon, namely a complete
bipartite geometry. A vertex of the second type has $[B:H]$ neighbours
among $[G:A]$ vertices of the first type. By transitivity, completeness
is therefore equivalent to equality of these numbers. The two alternatives
are exactly the required intersection property.
\end{proof}

\section*{Application to $\operatorname{PSL}(3,4)$}
The computation uses the degree-$21$ permutation representation supplied
by GAP. It acquires the complete maximal-subgroup algorithms and data for
$G$, its maximal subgroups, and the intermediate groups encountered during
the overgroup search. These subgroup classifications are retained inputs;
the argument above does not reprove them. In particular, the routine for
intermediate subgroups must return actual subgroups, rather than only
conjugacy representatives.

There are nine conjugacy classes of maximal subgroups. Their local lists
give $42$ pairs $(A,H)$. The third step yields $60$ pairs $(A,B)$, of which
$31$ pass the incidence test. They give the following seven classes.
The source labels refer to the complete Halle preprint of Gottschalk and
Leemans, \emph{Geometries for the group PSL(3,4)}, Report 99--30 (1999),
pp.~7--8; its later journal edition has not been collated here.
\[
\begin{array}{crrrrr}
\toprule
\text{Source}&|A|&|B|&|H|&[A:H]&[B:H]\\
\midrule
2.1&960&960&192&5&5\\
2.2&360&96&24&15&4\\
2.3&360&72&36&10&2\\
2.4&168&168&21&8&8\\
2.5&168&96&24&7&4\\
2.6&72&72&8&9&9\\
2.7&72&16&8&9&2\\
\bottomrule
\end{array}
\]
All surviving local actions are doubly transitive. This is an output of
the computation: the search includes primitive actions that are not
doubly transitive and rejects pairs only by the stated tests. The seven
unordered pairs of parabolic orders are distinct. Since the two type
sizes are $[G:A]$ and $[G:B]$, the seven representatives are also
inequivalent as abstract incidence geometries, even allowing interchange
of types. Explicit automorphisms identifying the other surviving pairs
are supplied separately as images of generators.

The coset correspondence and the stabilizer criterion for primitivity
are classical. The present argument organizes these inputs into the same
finite construction used for doubly transitive residues, with the
incidence test extended to all local pairs. No claim of historical
novelty for this enumeration principle is made.
\end{document}
